\section{Davenport-Hasse Theorem}
So using the discussion on \emph{Dirichlet $L$-functions} we can relate Gaussian sums in $\bbF_q$ with Gaussian sum of lifted characters to any finite extension of $\bbF_q$. 
\begin{Theorem}{Davenport-Hasse Theorem}{thm:davenport-hasse}
  Let $\chi\in\Xq$ be an additive character and $\psi\in\Mq$ be a multiplicative character of $\bbF_q$, not both of them trivial. Suppose $\chi$ and $\psi$ are lifted to character $\chi^{(r)}$ and $\psi^{(r)}$ respectively of the finite extension field $E$ of $\bbF_q$ with $[E\colon \bbF_q]=r$. Then $$G(\psi^{(r)},\chi^{(r)})=(-1)^{r-1}G(\psi,\chi)^r$$
\end{Theorem}
\begin{proof}
  Like previous section $G$ is the set of all rational functions $r(X)=\nicefrac{h_1}{h_2}$ where $h_1,h_2\in\Phi$, $h_2$ non-zero. Let $\ovG$ be the subgroup of $G$ consisting of rational functions $r(X)=\frac{h_1(X)}{h_2(X)}$ having $h_1(0),h_2(0)\neq 0$. Then define the multiplicative function $\lm:G\to \bbC$ such that for all $r\in G$, $|\lm(r)|\leq 1$ and $\lm(1)=1$ in the following way: 

  If $r(X)\in G$ where $r(X)=\frac{h_1(X)}{h_2(X)}$ and in the common splitting field $E$ of $h_1,h_2$, they factor as: $$h_1(X)=\prod_{i=1}^{d_1}(X-\alpha_i)\qquad \text{and}\qquad h_1(X)=\prod_{j=1}^{d_2}(X-\beta_j)$$ where $\alpha_i,\beta_j\in E$ for all $i\in[d_1]$, $j\in[d_2]$ then $$\lm(r)=\psi\lt(\prod_{i=1}^{d_1}\alpha_i \cdot \prod_{j=1}^{d_2}\beta_j^{-1}\rt)\cdot \chi\lt(\sum_{i=1}^{d_1}\alpha_i-\sum_{j=1}^{d_2}\beta_j\rt)$$Now for any monic polynomial $h\in\Phi$ let $h(X)=X^n-c_1\cdot X^{n-1}+\cdots +(-1)^nc_n$ where $c_i\in\bbF_q$ for all $i\in[n]$. Then we have $\lm(h)=\psi(c_n)\cdot \chi(c_1)$. Clearly $\lm$ is a multiplicative map on $G$. Now for a given $(c_1,c_k)$ pair there are $q^{n-2}$ polynomials where same $(c_1,c_n)$ occurs. Hence $$\sum_{h\in\Phi_k}\lm(h)=q^{n-2}\sum_{c_1\in\bbF_q}\sum_{c_n\in\bbF_q^*}\psi(c_n)\cdot \chi(c_1)=q^{n-2}\lt(\sum_{c_n\in\bbF_q}\rt)\lt(\sum_{c_1\in\bbF_q}\chi(c_1)\rt)$$Since $\psi$ and $\chi$ are non-trivial by \thmref{thm:add-char-props} and \thmref{thm:mult-char-props} we have $\sum_{h\in\Phi_k}\lm(h)=0$ for $n>1$. 

  Now $\Phi(0)=\{1\}$ and $\Phi(1)=\{x-\alpha\colon \alpha\in\bbF_q^*\}$. Hence $$\sum_{h\in\Phi_1}\lm(h)=\sum_{\alpha\in\bbF_q^*}\psi(\alpha)\cdot \chi(\alpha)=G(\psi,\chi)$$Therefore $$\ovL(z,\lm)=1+G(\psi,\chi)\cdot z\qquad\text{and}\qquad L(s,\chi)=1+q^{-s}\cdot G(\psi,\chi)$$Therefore root of $\ovL(z,\lm)$ is $-G(\psi,\chi)$.  

  Now we will calculate $L_r$  where $r=[E\colon \bbF_q]$. From the definition of $L_r$ in previous section we have $$L_r =  \sum_{\substack{h\in\irr(\Phi),\\ \deg(h)\mid r}}\deg(h)\cdot \chi(h)^{\nicefrac{r}{\deg(h)}}= \sum_{\substack{h\in\irr(\Phi),\\ \deg(h)\mid r}} \deg(h)\cdot \lm\lt( h^{\nicefrac{r}{\deg(h)}}\rt)$$Now for any irreducible $h$ let $\gm$ be a root of $h$ in the splitting field of $f$. Let $$h^{\nicefrac{r}{\deg(f)}}= x^r-c_1\cdot x^{r-1}+\cdots+(-1)^k c_r=\prod_{i=o}^{r-1}\lt(X-\gm^{q^i}\rt)$$Hence 
  $c_1=\tr_{E/\bbF_q}(\gm)$ and $c_r=\norm_{E/\bbF_q}(\gm)$. Hence we have $$\lm\lt(h^{\nicefrac{r}{\deg(h)}}\rt)=\psi(c_r)\cdot \chi(c_1)=\psi\lt(\norm_{E/\bbF_q}(\gm)\rt)\cdot \chi\lt(\tr_{E/\bbF_q}(\gm)\rt)=\psi^{(r)}(\gm)\cdot \chi^{(r)}(\gm)$$Now we can express $\deg(h)\cdot \lm\lt(h^{\nicefrac{r}{\deg(h)}}\rt)=\sum_{\gm\in E,h(\gm)=0}\psi^{(r)}(\gm)\cdot \chi^{(r)}(\gm)$. 
  Therefore $$L_r= \sum_{\substack{h\in\irr(\Phi),\\ \deg(h)\mid k}} \deg(h)\cdot \lm\lt( h^{\nicefrac{k}{\deg(h)}}\rt)=\sum_{\substack{h\in\irr(\Phi),\\ \deg(h)\mid k}}\sum_{\gm\in E,h(\gm)=0}\psi^{(r)}(\gm)\cdot \chi^{(r)}(\gm) $$Now all irreducible polynomials of $\bbF_q$ whose degree divides $r$ with non-zero constant their roots in $E$ are disjoint sets and partition $E^*$. Therefore we have $$L_r=\sum_{\gm\in E^*}\psi^{(r)}(\gm)\cdot \chi^{(r)}(\gm)=G(\psi^{(r)},\chi^{(r)})$$. Since  by \thmref{thm:l-function-root-factorisation} we have $L_r=-(-G(\psi,\chi))^r$. So we have the theorem. 
\end{proof}

We will see extension of this process of defining $\lm$ for special equations over finite fields in chapter 7.

For certain special characters, the associated Gaussian sums can be evaluated explicitly. A very special multiplicative character to study is the quadratic character, $\eta$. We get a very celebrated formula for the Gaussian sum of $\eta$ and the canonical additive character. 
\begin{Theorem}{}{}
  Let $\bbF_q$ be a finite field with $q=p^r$ where $p$ is an odd prime and $r\in\bbN$. Let $\eta$ be the quadratic character of $\bbF_q$ and let $\chi_1$ be the canonical additive character of $\bbF_q$. Then $$G(\eta,\chi_1)=\begin{cases}
    (-1)^{r-1}\cdot q^{\nicefrac12} & \text{if }p\equiv 1\bmod 4\\ 
    (-1)^{r-1}\cdot i^{r}\cdot q^{\nicefrac12} & \text{if }p\equiv 3\bmod 4    
  \end{cases}$$
  In other words, $$G(\eta,\chi_1)=(-1)^{r-1}\cdot i^{\frac{(p-1)^2}4 r} \cdot q^{\nicefrac12}$$ 
\end{Theorem}
\begin{proof}
  Since $\eta$ always takes values in $\pm 1$ we have $\eta=\ov{\eta}$. Hence by using \lemref[(iv)]{lem:gaussian-sum-identity} we have $G(\eta,\chi_1)^2=\eta(-1)q$. Now if $q\equiv 1\bmod 4$ then $\eta(-1)=1$ and if $q\equiv 3\bmod 4$ then $\eta(-1)=-1$. Hence \begin{equation}
    G(\eta,\chi_1)=\begin{cases}
      \pm\ q^{\nicefrac12}& \text{if }q\equiv 1\bmod 4\\ 
      \pm\ i\cdot q^{\nicefrac12}& \text{if }q\equiv 3\bmod 4
    \end{cases}
  \end{equation}
  Since $i^{\nicefrac{(p-1)^2}{4}}=\begin{cases}
      1 & \text{if }q\equiv 1\bmod 4\\ 0 & \text{if }q\equiv 3\bmod 4
    \end{cases}$

  \begin{observation}\label{obs:gauss-quad-char}
    For any finite field $\bbF_q$,  $G(\eta,\chi_1)=\pm\ i^{\nicefrac{(q-1)^2}{4}}\cdot q^{\nicefrac12}$ 
  \end{observation}

  Now comes the hardest part of the proof that is getting the signs correctly. We will first prove this for $r=1$ case then we will go for $r>1$. Now we have the observation
  
  
  Let $W$ be the vector space of all complex valued functions on $\bbF_p^*$. It is a $(p-1)$-dimensional vector space over $\bbC$. So $\mcB\coloneqq \{f_1,\dots, f_{p-1}\}$ where $f_j(k)=1$ if $j=k$ and other wise $0$ forms a basis of $W$ for all $j\in[p-1]$. Now by using the orthogonal properties of multiplicative characters in \thmref{thm:mult-char-props} the set $\mcB_{\Mq}\coloneqq \{\psi_0,\dots, \psi_{p-2}\}$ as defined in \lemref{lem:mult-char-define}, also forms a basis of $W$.

  Let $\zeta_p$ is the $p^{th}$ root of unity i.e. $\zeta_p=e^{\nicefrac{2\pi i}{p}}$. Hence if $\chi_k\in\Xp$ be an additive character of $\bbF_p$ as defined in \thmref{thm:add-char-define}, then for all $j\in\bbF_p$, $\chi_k(j)=\zeta_p^{jk}$. Then define the linear operator $T:W\to W$ in the following way:
  \begin{equation}\label{eq:quad-linear-operator}
    \forall\ h\in W,\ \forall\ k\in[p-1]\qquad (Th)(k)=\sum_{j=1}^{p-1}\zeta_p^{jk}h(j)
  \end{equation}
  Therefore for all basis elements $\psi_i$ $$(T\psi_i)(k)=\sum_{j=1}^{p-1}\zeta_p^{jk}\cdot \psi_i(k)=\sum_{j=1}^{p-1}\chi_k(j)\cdot \psi_i(k)=G(\psi,\chi_k)=\ov{\psi_i}(k)\cdot G(\psi_i,\chi_1)\quad \forall\ k\in\bbF_p^*$$where the last equality comes from \lemref[(i)]{lem:gaussian-sum-identity}. Hence $T\psi_i=G(\psi_i,\chi_1)\cdot \ov{\psi_i}$ for every multiplicative characters of $\bbF_p$. 

  Since $\ov{\psi}=\psi$ only for the trivial character and quadratic character, the matrix for $T$ in the basis $\mcB_{\Mq}=\{\psi_0,\dots, \psi_{p-2}\}$ contains two diagonal entries, $G(\psi_0,\chi_1)=-1$ and $G(\eta,\chi_1)$ for $\psi_0$ and $\eta=\psi_{\nicefrac{p-1}2}$. For all other characters we can divide them into groups of pairs where the characters in each pair are conjugates of each other. Then for each pair $(\psi,\ov{\psi})$ we have the following block matrix in the matrix of $T$. $$\mat{0 & G(\ov{\psi},\chi_1)\\ G(\psi,\chi_1)& 0}$$Therefore in the determinant of the matrix $T$ each block contributes $-G(\psi,\chi_1)G(\ov{\psi},\chi_1)=-\psi(-1)p$ by \lemref[(iv)]{lem:gaussian-sum-identity}. So we get $$\det(T)=(-1)\cdot G(\eta,\chi_1) \cdot (-p)^{\nicefrac{(p-3)}2}\prod_{j=1}^{\frac{p-3}2}\psi_j(-1)$$Let $\psi\in\Mp$ is the principle multiplicative character of $\bbF_p$. Then $\psi_j=\psi^j$. Therefore $\psi_j(-1)=\psi^j(-1)=(-1)^j$ as $\psi(-1)=-1$. Hence we have \begin{align*}
    \det(T) & = - G(\eta,\chi_1) \cdot (-p)^{\nicefrac{(p-3)}2}\prod_{j=1}^{\frac{p-3}2}(-1)^j\\ 
    & = - G(\eta,\chi_1) \cdot (-p)^{\nicefrac{(p-3)}2}\cdot (-1)^{\sum\limits_{j=1}^{\nicefrac{(p-3)}{2}}j}\\ 
    & =  (-1)^{\nicefrac{(p-1)}{2}}\cdot G(\eta,\chi_1) \cdot p^{\nicefrac{(p-3)}2}\cdot (-1)^{\nicefrac{(p-1)(p-3)}8}\\ 
    & = (-1)^{\nicefrac{(p-1)}{2}}\cdot \lt(\pm\ i^{\nicefrac{(p-1)^2)}{4}}\cdot p^{\nicefrac12}\rt) \cdot p^{\nicefrac{(p-3)}2}\cdot (-1)^{\nicefrac{(p-1)(p-3)}8}\by{\obsref{obs:gauss-quad-char}}\\ 
    & = \pm (-1)^{\nicefrac{(p-1)}{2}}\cdot i^{\frac{(p-1)^2)}{4}+\frac{(p-1)(p-3)}4}\cdot p^{\nicefrac{(p-2)}{2}}\\ 
    & = \pm (-1)^{\nicefrac{(p-1)}{2}}\cdot i^{\nicefrac{(p-1)(p-2)}2}\cdot p^{\nicefrac{(p-2)}{2}}
  \end{align*}
  So we get \begin{equation}\label{eq:fourier-basis-det}
    \det(T)=\pm (-1)^{\nicefrac{(p-1)}{2}}\cdot i^{\nicefrac{(p-1)(p-2)}2}\cdot p^{\nicefrac{(p-2)}{2}}
  \end{equation}

  Now we will compute $\det(T)$ utilizing the basis $\mcB=\{f_1,\dots, f_{p-1}\}$. From the construction of $T$ in \eqref{eq:quad-linear-operator} we get the $(j,k)^{th}$ element of the matrix of $T$ in the $\{f_1,\dots, f_{p-1}\}$ basis to be $\zeta_p^{jk}$ for any $0\leq j,k\leq p-1$. So we have \begin{align*}
    \det(T)& =\det\lt( \lt\{\zeta_p^{jk}\rt\}_{1\leq j,k\leq p-1}\rt)\\ 
    & = \det\lt( \lt\{\zeta_p^j\zeta_p^{j(k-1)}\rt\}_{1\leq j,k\leq p-1}\rt)\\ 
    & = \zeta_p^{1+2+\cdots+(p-1)}\det\lt( \lt\{\zeta_p^{j(k-1)}\rt\}_{1\leq j,k\leq p-1}\rt)\\ 
    & = \det\lt( \lt\{\zeta_p^{j(k-1)}\rt\}_{1\leq j,k\leq p-1}\rt) & \lt[\text{As }p\mid \frac{p(p-1)}2=1+2+\cdots+(p-1)\rt]
  \end{align*}
  Now the matrix $\lt\{\zeta_p^{j(k-1)}\rt\}_{1\leq j,k\leq p-1}$ is the \emph{Vandermonde} matrix $V(\zeta_p,\zeta_p^2,\dots, \zeta_p^{p-1})$. Therefore $$\det(T)=\prod_{1\leq j<k\leq p-1}\lt(\zeta_p^k-\zeta_p^j\rt)$$We will show that 
  \begin{equation}\label{eq:quadratic-gaussian-det}
    (-1)^{\nicefrac{(p-1)}{2}}\cdot i^{\nicefrac{(p-1)(p-2)}2}\cdot p^{\nicefrac{(p-2)}{2}}=\prod_{1\leq j<k\leq p-1}\lt(\zeta_p^k-\zeta_p^j\rt)
  \end{equation}
  So let $\zeta_{2p}=e^{\nicefrac{\pi i}{p}}$. Hence $\zeta_{2p}^2=\zeta_p$. By the properties of complex numbers we know for all $t\in\bbZ$, $\zeta_{2p}^t=\cos\lt(\frac{t\pi}{p}\rt)+i\sin\lt(\frac{t\pi}p\rt)$. Hence $\zeta_{2p}^t-\zeta_{2p}^{-t}=2i\sin\lt(\frac{t\pi}p\rt)$. So we get \begin{align*}
    \det(T) & = \prod_{1\leq j<k\leq p-1}\lt(\zeta_p^k-\zeta_p^j\rt)\\
    & = \prod_{1\leq j<k\leq p-1}\lt(\zeta_{2p}^{2k}-\zeta_{2p}^{2j}\rt)\\
    & = \prod_{1\leq j<k\leq p-1}\zeta_{2p}^{j+k}\lt(\zeta_{2p}^{k-j}-\zeta_{2p}^{-(k-j)}\rt)\\ 
    & =  \lt[\prod_{1\leq j<k\leq p-1}\zeta_{2p}^{j+k} \rt]\;\lt[\prod_{1\leq j<k\leq p-1} \lt(2i\sin\frac{\pi(k-j)}{p}\rt) \rt]\\ 
    & = \zeta_{2p}^{\sum_{1\leq j<k\leq p-1}(j+k)} \cdot i\st^{\binom{p-1}{2}}\prod_{1\leq j<k\leq p-1} \lt(2\sin\frac{\pi(k-j)}{p}\rt) \\ 
    & = \zeta_{2p}^{\sum_{1\leq j<k\leq p-1}(j+k)} \cdot i{\nicefrac{(p-1)(p-2)}  {2}}\prod_{1\leq j<k\leq p-1} \lt(2\sin\frac{\pi(k-j)}{p}\rt) 
  \end{align*}
  Now we calculate the power of $\zeta_{2p}$ in the first term. $$\sum_{1\leq j<k\leq p-1}(j+k)=\sum_{k=2}^{p-1}\sum_{j=1}^{k-1}(j+k)=\sum_{k=2}^{p-1}\lt(\frac{k(k-1)}2+k(k-1)\rt)=\frac32\sum_{k=2}^{p-1}(k^2+k)=\frac{p(p-1)(p-2)}2$$Hence $$\zeta_{2p}^{\sum_{1\leq j<k\leq p-1}(j+k)} = \lt(\zeta_{2p}^p\rt)^{\nicefrac{(p-1)(p-2)}{2}}=(-1)^{\nicefrac{(p-1)(p-2)}{2}}=\lt((-1)^{p-2}\rt)^{\nicefrac{(p-1)}{2}}=(-1)^{\nicefrac{(p-1)}{2}}$$So we have \begin{equation}\label{eq:standard-basis-det}
    \det(T)=(-1)^{\nicefrac{(p-1)}{2}}\cdot i^{\nicefrac{(p-1)(p-2)}{2}}\cdot \prod_{1\leq j<k\leq p-1} \lt(2\sin\frac{\pi(k-j)}{p}\rt)
  \end{equation}Notice if we compare this with \eqref{eq:fourier-basis-det}, the first two terms are already same. If we can show that $\prod_{1\leq j<k\leq p-1} 2\sin\frac{\pi(k-j)}{p}>0$ then we have \eqref{eq:quadratic-gaussian-det}. 
  
  Now since $j,k\in[p-1]$ and $k>j$ we have $k-j\in[p-1]$. Hence $\sin\frac{\pi(k-j)}{p}>0$ for all $1\leq j<k\leq p-1$. Hence $$\prod_{1\leq j<k\leq p-1} \lt(2\sin\frac{\pi(k-j)}{p}\rt) >0$$ So if we compare \eqref{eq:standard-basis-det} with \eqref{eq:fourier-basis-det} and match the signs we get $$\det(T)=(-1)^{\nicefrac{(p-1)}{2}}\cdot i^{\nicefrac{(p-1)(p-2)}2}\cdot p^{\nicefrac{(p-2)}{2}}$$So from this we get $G(\eta,\chi_1)=i^{\nicefrac{(p-1)^2}{4}}\cdot p^{\nicefrac12}$ since at no point of our calculation the sign was changed. 

  So now we have established the theorem for $r=1$. Now for $r>1$ following \autoref{sec:lifting-chars} we can lift the quadratic character of $\bbF_p$, $\eta$ to a multiplicative character $\eta^{(r)}$ of $\bbF_q$ where $q=p^r$ and $\eta^{(r)}=\eta\circ N$ where $N$ is the absolute norm function from $\bbF_q$ to $\bbF_p$. Since $\eta$ is a quadratic character $\eta^{(r)}$ will be the quadratic character of $\bbF_q$. Similarly we can lift the canonical additive character $\chi_1$ of $\bbF_p$ to the canonical additive character $\chi_q^{(r)}$ of $\bbF_q$ by $\chi_1^{(r)}=\chi_1\circ \tr$ where $\tr$ is the absolute trace function from $\bbF_q$ to $\bbF_p$. Then by using \DHthm we get $$    G(\eta^{(r)},\chi_1^{(r)}) = (-1)^{r-1}G(\eta,\chi_1)^r = (-1)^{r-1}\lt[i^{\nicefrac{(p-1)^2}{4}}\cdot p^{\nicefrac12}\rt]=(-1)^r \cdot i^{\frac{(p-1)^2}{4}r}\cdot p^{\nicefrac{r}2}=(-1)^r \cdot i^{\frac{(p-1)^2}{4}r}\cdot q^{\nicefrac12}$$So we have the theorem.
\end{proof}

The proof of the above theorem gives us another important result. 
\begin{lemma}{}{}
  If $p$ is a odd prime and $\zeta_p$ is the $p^{th}$ root of unity. Let $V(\zeta_p,\zeta_p^2,\dots,\zeta_p^{p-1})$ is the $p\times p$ Vandermonde Matrix. Then $$\det\lt(V(\zeta_p,\zeta_p^2,\dots,\zeta_p^{p-1})\rt)=(-1)^{\nicefrac{(p-1)}{2}}\cdot i^{\nicefrac{(p-1)(p-2)}2}\cdot p^{\nicefrac{(p-2)}{2}}$$
\end{lemma}